\subsection{Position Vectors and Displacement Vectors}

Let the origin of the coordinate system be $O=(0,0)$, and let
\[
A=(x_1,y_1),
\qquad
B=(x_2,y_2).
\]

The arrow from the origin to a point records the position of that point relative
to the chosen origin. We write
\[
\boxed{\mathbf r_A=\overrightarrow{OA}=(x_1,y_1)},
\qquad
\boxed{\mathbf r_B=\overrightarrow{OB}=(x_2,y_2)}.
\]
These are the \emph{position vectors} of $A$ and $B$.

If an object moves from $A$ to $B$, we are no longer asking where either point
is. We are asking what change carries the initial position to the final one. The
answer is the \emph{displacement vector}
\[
\Delta\mathbf r=\overrightarrow{AB}.
\]
Its components are obtained by subtracting the initial coordinates from the
final coordinates:
\[
\boxed{
\overrightarrow{AB}
=(x_2-x_1,\,y_2-y_1)
}.
\]

The relation between the two kinds of vectors is therefore
\[
\boxed{
\Delta\mathbf r
=\overrightarrow{AB}
=\mathbf r_B-\mathbf r_A
}.
\]
Equivalently,
\[
\boxed{
\mathbf r_B=\mathbf r_A+\Delta\mathbf r
}.
\]
The final position is obtained by adding the displacement to the initial
position.

\begin{center}
\begin{tikzpicture}[scale=0.9]
    \draw[axis] (-0.5,0) -- (7.0,0) node[right] {$x$};
    \draw[axis] (0,-0.5) -- (0,5.2) node[above] {$y$};

    \coordinate (O) at (0,0);
    \coordinate (A) at (2,1.4);
    \coordinate (B) at (5.4,4.0);

    \draw[helper,->] (O) -- (A)
        node[midway,below right] {$\mathbf r_A$};
    \draw[helper,->] (O) -- (B)
        node[midway,above left] {$\mathbf r_B$};
    \draw[subset,->] (A) -- (B)
        node[midway,below right] {$\Delta\mathbf r=\overrightarrow{AB}$};

    \fill (O) circle (2pt) node[below left] {$O$};
    \fill[geometricSubsetColor] (A) circle (2pt) node[below right] {$A$};
    \fill[geometricSubsetColor] (B) circle (2pt) node[above right] {$B$};
\end{tikzpicture}
\end{center}

A position vector depends on the chosen origin: moving the origin changes the
coordinates assigned to the same point. A displacement vector compares two
positions. If the whole coordinate system is translated, both position vectors
change in the same way, but their difference remains the same.

The same vector can also be drawn from another starting point. If
$P=(p_1,p_2)$ and
\[
\mathbf v=(v_1,v_2),
\]
then applying this displacement at $P$ produces
\[
Q=P+\mathbf v=(p_1+v_1,\,p_2+v_2).
\]
Thus a displacement vector is not tied to one particular place: it describes
the same change wherever it is applied.

The zero vector
\[
\boxed{\mathbf 0=(0,0)}
\]
means no displacement. The opposite vector
\[
\boxed{-\mathbf v=(-v_1,-v_2)}
\]
reverses a displacement. In particular,
\[
\boxed{\overrightarrow{BA}=-\overrightarrow{AB}}.
\]

For motion depending on time, the position will later be written as
$\mathbf r(t)$. The displacement between two moments $t_1$ and $t_2$ will then
have exactly the same form:
\[
\Delta\mathbf r
=\mathbf r(t_2)-\mathbf r(t_1).
\]
